% MTYM — Note de synthèse
% Modèle adapté du gabarit officiel TFJM², à compléter par les équipes
% Opposant et Rapporteur de chaque passage.

\documentclass[11pt,a4paper]{article}

\usepackage[utf8]{inputenc}
\usepackage[T1]{fontenc}
\usepackage[french]{babel}
\usepackage[margin=2.2cm]{geometry}
\usepackage{amsmath,amssymb}
\usepackage{enumitem}
\usepackage{array}
\usepackage{titlesec}
\usepackage{parskip}

\setlist{nosep,leftmargin=1.4em}
\titleformat{\section}{\normalfont\Large\bfseries}{}{0pt}{}
\titlespacing*{\section}{0pt}{*1.6}{*0.8}

\newcommand{\blank}[1]{\rule{#1}{0.4pt}}
\newcommand{\smallblank}{\blank{2.2em}}
\newcommand{\midblank}{\blank{6em}}

% Table cell helper — 8 numbered rows
\newcommand{\synthtable}{%
  \begin{tabular}{|c|c|c|c|c|}
  \hline
  Question & ER & PR & QE & NT \\\hline
  \phantom{Q1} & & & & \\\hline
  \phantom{Q1} & & & & \\\hline
  \phantom{Q1} & & & & \\\hline
  \phantom{Q1} & & & & \\\hline
  \phantom{Q1} & & & & \\\hline
  \phantom{Q1} & & & & \\\hline
  \phantom{Q1} & & & & \\\hline
  \phantom{Q1} & & & & \\\hline
  \end{tabular}%
}

\begin{document}

\begin{center}
  {\Huge\bfseries MTYM}\\[0.3em]
  {\Large NOTE DE SYNTHÈSE}
\end{center}

\vspace{0.8em}
Poule \smallblank\\
Problème \smallblank{} défendu par l'équipe \midblank\\
Synthèse par l'équipe \midblank{} dans le rôle de :
$\square$ Opposant \quad $\square$ Rapporteur

\section*{Évaluation question par question de la solution}

\noindent
\begin{minipage}{0.40\textwidth}\synthtable\end{minipage}\hfill
\begin{minipage}{0.40\textwidth}\synthtable\end{minipage}

\medskip

\noindent
\textbf{ER} : entièrement résolue, ni erreur, ni manque mathématique\\
\textbf{PR} : partiellement résolue\\
\textbf{QE} : quelques éléments de réponse\\
\textbf{NT} : non traitée

\smallskip
\emph{Remarque :} il est possible de cocher entre les cases pour un cas
intermédiaire.

\section*{Erreurs et imprécisions}

Listez ci-dessous par ordre décroissant d'importance au plus quatre
erreurs et/ou imprécisions selon vous, en précisant la question
concernée, la page, le paragraphe et le type de remarque.

\vspace{0.4em}
\newcommand{\erreurEntry}[1]{%
  #1. Question \smallblank{} Page \smallblank{} Paragraphe \smallblank\\
  $\square$ Erreur majeure \quad $\square$ Erreur mineure \quad
  $\square$ Imprécision \quad $\square$ Autre : \blank{4em}\\
  Description :\\[4em]
}
\erreurEntry{1}
\erreurEntry{2}
\erreurEntry{3}
\erreurEntry{4}

\section*{Aspects positifs}

Identifiez au plus deux points forts spécifiques de la solution et
dites pourquoi (exemples : propositions majeures, idées importantes,
généralisations pertinentes, exemples significatifs, constructions
originales,\ldots).

\vspace{0.4em}
\newcommand{\positifEntry}[1]{%
  #1. Question \smallblank{} Page \smallblank{} Paragraphe \smallblank\\
  Description :\\[4em]
}
\positifEntry{1}
\positifEntry{2}

\section*{Évaluation qualitative de la solution}

Donnez votre avis concernant la solution en général. Mettez notamment
en valeur ses qualités globales, et précisez ce qui aurait pu
l'améliorer.

\vfill

\begin{center}
\textbf{Évaluation générale :} \quad
$\square$ Excellente \quad $\square$ Bonne \quad
$\square$ Suffisante \quad $\square$ Passable
\end{center}

\section*{Autres remarques (facultatif)}

Présentation, lisibilité, orthographe, etc.

\vfill
\hrulefill\\
\small Modèle rédigé par TFJM\textsuperscript{2}, adapté pour MTYM.

\end{document}
